\documentclass[12pt]{article}
\usepackage[utf8]{inputenc}
\usepackage[T1]{fontenc}
\usepackage{hyperref}
\title{Mazatl\'an Mazatec Language Resources for K-12 Curriculum Development}
\author{Compiled by Claude Sonnet 5 (claude-sonnet-5) for Andr\'as Kornai}
\date{}
\begin{document}
\maketitle

\section{Introduction}
Mazatl\'an Mazatec (ISO 639-3 vmz; Glottocode maza1296) is an Otomanguean language spoken in Oaxaca, Mexico, around Mazatl\'an Villa de Flores, with about 12,000 speakers. The search found Scripture and audio materials but no dictionary, grammar or school materials.

\section{Primary Educational Resources}
No primers, textbooks or school programmes were found in the search results.

\section{Linguistic and Reference Materials}
\begin{itemize}
\item HLT language database entry for vmz. It lists no online lexical resources, and its counts of online descriptions do not identify specific works: \url{https://hlt.bme.hu/en/dld/language/1776}
\item Wikipedia entry on the language, with speaker figures: \url{https://www.wikipedia.com/wiki/ISO_639:vmz}
\end{itemize}

\section{Religious and Community Resources}
\begin{itemize}
\item Scripture Earth page for vmz, with a readable New Testament book and downloadable New Testament audio: \url{https://www.scriptureearth.org/00eng.php?iso=vmz}
\end{itemize}

\section{Digital and Audio Resources}
\begin{itemize}
\item Global Recordings Network: audio and slideshow lessons, including the ``Good News'' set of 40 audio-visual Bible lessons: \url{https://globalrecordings.net/en/language/vmz}
\end{itemize}

\section{Curriculum Development Context}
The language is stable and used by all in the community, but no school use is known. Very little documentation was found, so a curriculum would depend on new material. Searches for a Spanish-titled dictionary, SIL materials and Ethnologue would be the next step. The Scripture and audio recordings offer text and speech samples.
\end{document}